\documentclass[fleqn]{unifal}          % classe base para a monografia

%==============================================================================
% Utilizacao de pacotes
\setlength{\headheight}{14.49998pt} %altura do cabeçalho
\setlength{\footskip}{4.35004pt}    %altura do rodape
\usepackage[T1]{fontenc}         % usa fontes postscript com acentos
\usepackage[brazilian]{babel}    % hifenização e títulos em português do Brasil
\usepackage[utf8]{inputenc}     % permite edição direta com acentos
\usepackage[fleqn]{amsmath}             % pacote da AMS para Matemática Avançada
\usepackage{amssymb}             % símbolos extras da AMS
\setlength{\mathindent}{1.25cm}   % alinhamento das equações
\usepackage{latexsym}            % símbolos extras do LaTeX
\usepackage{ragged2e}            % alinhamento adicional
\usepackage{graphicx}            % para inserção de gráficos
\usepackage{tikz}                % figuras vetoriais
\usepackage{pgfplots}            % gráficos vetoriais
\usepackage{listings}            % para inserção de código
\usepackage{fancyvrb}            % para inserção de saídas de comandos
%\usepackage{enumerate}           % para personalizar lista enumeradas 
%(incluso na classe)
\usepackage{longtable}           % para tambelas muito grandes NOVO!!!!

\usepackage{colortbl} % cores em tabelas
\newcolumntype{Z}{|>{\columncolor[gray]{0.9}}l|} %cor cinza em células
%\usepackage{array} % já incluso na classe
\newcolumntype{L}[1]{>{\raggedright\let\newline\\\arraybackslash\hspace{0pt}}m{#1}}
\newcolumntype{C}[1]{>{\centering\let\newline\\\arraybackslash\hspace{0pt}}m{#1}}
\newcolumntype{R}[1]{>{\raggedleft\let\newline\\\arraybackslash\hspace{0pt}}m{#1}}
\usepackage{multirow} % para juntar duas linhas em uma só

\usepackage{multicol} % para uso de várias colunas

% cores para os links cruzados
\usepackage{color}
\definecolor{rltred}{rgb}{0.2,0,0}
\definecolor{rltgreen}{rgb}{0,0.2,0}
\definecolor{rltblue}{rgb}{0,0,0.2}

\usepackage[colorlinks=true,
urlcolor=rltblue,       % \href{...}{...} external (URL)
filecolor=rltgreen,     % \href{...} local file
linkcolor=rltred,       % \ref{...} and \pageref{...}
citecolor=rltgreen,
pdftitle={Exemplo de Uso da Classe Uflamon},
pdfauthor={Joaquim Quinteiro Uchôa},
pdfsubject={Este texto tem por objetivo servir de exemplo da classe Uflamon.},
pdfkeywords={Comunicação Científica. 2. Pesquisa . 3. Pesquisa Científica. 
	4. Redação. 5. Monografia.}%
]{hyperref} % para referência cruzadas
\hypersetup{hypertexnames=false}
%\usepackage{hyperref}            % para referência cruzadas
\usepackage{subfigure}           % figuras dentro de figuras
\usepackage{caption}            % remodelando o formato dos títulos de 
% tabelas e figuras
\usetikzlibrary{arrows.meta,positioning,shapes.geometric,fit,mindmap,calc,decorations.pathmorphing}
\pgfplotsset{compat=1.18}

% configuração padrão do listings   
\lstset{
	language=Java,
	extendedchars=true,
	tabsize=3,
	basicstyle=\footnotesize\ttfamily,
	stringstyle=\em,
	showstringspaces=false 
}

% fazendo equações com numeração sequencial
\counterwithout{equation}{chapter}

% para referências de acordo com a ABNT
% precisa instalar o abntex2 antes!!!
% http://abntex.codigolivre.org.br/
% comente se pretende usar outro padrão

%\usepackage[style=abnt]{biblatex}
%\renewcommand*{\mkbibnamefamily}[1]{#1}%
% arquivo de referência (adicione os arquivos de bibliografia, um por linha)
%\addbibresource{refbib.bib}
%abnt-emphasize=bf coloca o título das bibliografias em negrito
%abnt-thesis-year=both
\usepackage[alf,
abnt-nbr10520=1988,
abnt-etal-cite=3,
abnt-etal-list=3,
abnt-url-package=url,
abnt-emphasize=bf,
abnt-etal-text=it]{abntex2cite}
%######################################################################
% arrumando formato de referências com url no abntex
\def\UrlLeft{}
\def\UrlRight{}

%\usepackage[alf,abnt-repeated-title-omit=yes,abnt-emphasize=bf,abnt-etal-text = it,abnt-etal-list=0, abnt-url-package=url, abnt-full-initials=no]{abntex2cite}

% evite usar o hyperref com abntex, pode dar caca em urls... no linha anterior, informo
% para incluir urls usando o pacote url e não o hyperref
%
% caso queira o hyperref com abntex, comente a linha anterior e descomente a seguinte
%\usepackage[alf,abnt-etal-cite=3,abnt-etal-list=0,abnt-etal-text=emph]{abntex2cite}
%
% caso vc ainda use a versão anterior da abntex, comente a linha incluindo o abntex2cite
% e descomente a próxima linha 
%\usepackage[alf,abnt-etal-cite=3,abnt-etal-list=0,abnt-etal-text=emph]{abntcite}


% redefinindo formatação de títulos de tabelas e figuras


%==============================================================================
% para os fãs do Word, descomente as linhas abaixo
%\sloppy %mais espaço entre as linhas
%\usepackage{identfirst} %identando-se a primeira linha de cada seção
%\noindentfirst % Tire o comentário para manter o padrão do LaTeX.

%==============================================================================
% definido comandos na monografia - não é necessário na sua monografia 
% apenas para exemplificar a definição de novos comandos
\newcommand{\defs}[1]{\textsl{#1}}
\hfuzz=60pt
\vfuzz=40pt
\hbadness=10000
\vbadness=10000
\sloppy
\providecommand{\htmladdnormallink}[2]{\url{#2}}
\AtBeginDocument{\hfuzz=2pt\vfuzz=40pt}


% Especificando hifenizações que por ventura LaTeX não saiba fazer
% Por padrão 99,9% dos termos em português devem ser hifenizados corretamente.
\hyphenation{hardware software Li-nux am-bien-te diag-nos-ti-car coor-de-na-ção 
	FAE-PE Recovery TelEduc Williams UFLA}

%==============================================================================
% Dados da monografia, capa: autor, titulo, banca, etc... - SUBSTITUA DE ACORDO
%==============================================================================
\author{Rafael Passos Domingues (Grupo 3)\\
\small Espaço para inclusão dos demais integrantes: [Nome 2], [Nome 3], [Nome 4], [Nome 5], [Nome 6]}
\title{Modelo de Negócio para Startup de Base Tecnológica}
\subtitle{Integração entre Maturidade Tecnológica, Viabilidade Financeira e Incentivos Fiscais}
%\edicao{3$^a$ edição revista, atualizada e ampliada}
\date{2026}
\tipo{Trabalho acadêmico da disciplina de empreendedorismo e inovação do curso de Ciência da Computação da Universidade Federal de Alfenas.}
% use \orientador ou \orientadora quando for o caso
\orientador{Prof. Drd. Elivandro Miguel Reis}
%\orientadora{}
% use \coorientador ou \coorientadora quando for o caso
%\coorientadora{Profa. DSc. Maria Orientadora } % comente se não tiver coorientador
%\coorientador{}
\local{Alfenas/MG}
\bancaum{Prof. MSc. Antônio Banca Um}{UFM}
\bancadois{Prof. DSc. João Banca Dois}{FCO} % comente se sua banca tiver só um professor
\bancatres{Profa. Esp. Eliza Banca Três}{BELMIS}
\bancaquatro{Prof. Esp. Carlos Banca Quatro}{IBGPLUS}
\defesa{26 de março de 2026}
%==============================================================================
%##################################################
% Dados para Ficha catalográfica, gerada pelo sistema da Biblioteca da UFLA
% http://www.biblioteca.ufla.br/FichaCatalografica/
% dados para ficha catalográfica
% Elaboração da Ficha Catalográfica
%\preparofichacat{Sistema de Bibliotecas da Universidade Federal de Alfenas \\ Biblioteca Central}
% primeiro autor - como na primeira linha da ficha catalográfica
%\fcautor{Uchôa, Joaquim Quinteiro}
% autores, separados por vírgula - na ficha catalográfica, no formato que
% vem após o título e a barra ("/")
%\fcautores{Joaquim Quinteiro Uchôa}
% caso trabalho seja ilustrado (figuras, gráficos, tabelas, etc.), 
% então informar por meio do comando a seguir
% caso não seja ilustrado, basta comentá-lo
%\fcilustrado{il.}
% dados da edição para a ficha 
%\fcedicao{2$^a$ ed. rev., atual. e ampl.}
% tipo do trabalho (tese, dissertação, etc.), de acordo com sistema
% de geração de ficha catalográfica
%\fctipo{Tese(doutorado)}
% ano da defesa, só precisa informar se for diferente do ano da publicação
% se forem iguais, comente a linha a seguir
%\fcdatadefesa{2016}
% preencher aqui com os dados de catalogação gerados pelo sistema
%\fccatalogacao{1. TCC. 2. Monografia. 3. Dissertação. 4. Tese. 5. Trabalho Científico – Normas. I. Universidade Federal de Lavras. II. Título.}
%\fcclasi{808.066}

%##################################################

%\antesfichacat{\noindent Para citar este documento: \\UNIVERSIDADE FEDERAL DE LAVRAS. Biblioteca Universitária. \textbf{Manual de normalização e estrutura de trabalhos acadêmicos: TCC, monografias, dissertações e teses}. 2. ed. rev., atual. e ampl. Lavras, 2015. Disponível em: \url{http://www.biblioteca.ufla.br/wordpress/wpcontent/uploads/bdtd/manual_normalizacao_UFLA.pdf}. Acesso em: data de acesso.}

%\depoisfichacat{\noindent A reprodução e a divulgação total ou parcial deste trabalho são autorizadas, por qualquer meio convencional ou eletrônico, para fins de estudo e pesquisa, desde que citada a fonte.\\
	%\newline
	%{\small Este documento possui páginas em branco para facilitar a impressão frente-e-verso.}}

%##################################################

%##################################################

% para os exemplos do manual
%\newenvironment{exemplomanual}{
	%\vspace{0.5cm}
	%\noindent\begin{minipage}{\textwidth}
		%\noindent\rule{\textwidth}{0.5pt}
		%\vspace{-1cm}
		%\begin{flushleft}
		%}{
		%\end{flushleft}
		%\vspace{-0.6cm}
		%\noindent\rule{\textwidth}{0.5pt}
		%\vspace{0.3cm}
		%\end{minipage}
		%}
	
	%\newenvironment{exemplomanuallista}{
		%\vspace{0.3cm}
		%\noindent\begin{minipage}{\textwidth - 0.5cm}
			%\noindent\rule{\textwidth}{0.5pt}
			%\vspace{-1cm}
			%\begin{flushleft}
			%}{
			%\end{flushleft}
			%\vspace{-0.6cm}
			%\noindent\rule{\textwidth}{0.5pt}
			%\vspace{0.3cm}
			%\end{minipage}
			%}
		
		% por conta de alguns exemplos
		%\usepackage{setspace}
		
		%##################################################
		% se você gerou uma ficha catalográfica por algum sistema, basta
		% incluir o pdf ou jpg da ficha
		%\fichaCatalografica{fichaCatalografica}
		
		% se você precisar incluir alguma errata, por algum motivo,
		% basta incluir o pdf ou jpg da ficha
		%\errata{euerrei}
		
		% se vc já defendeu e tem o arquivo escaneado da folha de rosto, 
		% descomente e altere o nome do arquivo
		%\folhaAprovacaoAssinada{folhapro}
		
		% Aqui começa o documento propriamente dito
		\begin{document}
			
			\maketitle
			
			\dedic{A todo sócio de startup que precisou pivotar quando o plano ruiu.\\
\
Que a coragem de desfazer o que não sustenta seja sempre maior que o medo da perda.\\
\
Pois, às vezes, é preciso destruir o desenho antigo para construir uma estrutura mais resiliente, mais lúcida e mais viva.}     % Dedicatórias\\
			
			\thanks{Agradecemos ao Departamento de Ciência da Computação do ICEx da UNIFAL-MG pelo suporte acadêmico, técnico e institucional ao desenvolvimento deste trabalho.\\
				
				Agradecemos também ao orientador Prof. Drd. Elivandro Miguel Reis pelas contribuições críticas e pela condução metodológica ao longo do projeto.}         % Agradecimentos
			
			\epigrafe{ % citação opcional
				``Quem tem um porquê enfrenta qualquer como.''\\
				(Friedrich Nietzsche)}
			
			% palavras-chave
			\palchaves{modelo de negócio; startup de base tecnológica; business model canvas; gestão de valor; framework integrado.}
			\resumo{Este trabalho apresenta a estruturação de um modelo de negócio para startups de base tecnológica a partir de um framework integrado de cinco dimensões: estratégica, estrutural, engenharia, fiscal e governança de valor. A proposta organiza práticas de gestão, desenvolvimento e sustentabilidade financeira para reduzir incertezas na transição entre ideia e operação viável. O estudo destaca o uso de tailoring dinâmico, design de plataforma, métricas de fluxo de engenharia, aproveitamento da Lei do Bem e gestão por OKRs/OrKRs para apoiar decisões orientadas a valor. Como resultado, consolida-se um roteiro de implementação aplicável a projetos de inovação em Computação, com foco em validação de hipóteses, escalabilidade e consistência econômico-operacional.}  % Resumo (digite aqui o resumo)
			
			% keywords devem vir antes do abstract
			\keywords{business model; technology startup; platform design; value management; entrepreneurial strategy.} % keywords
			\abstract{This paper presents a business model structuring approach for technology-based startups using an integrated five-dimension framework: strategic, structural, engineering, fiscal and value governance. The proposal combines management, product development and financial sustainability practices to reduce uncertainty in the transition from idea to viable operation. The study emphasizes dynamic tailoring, platform design, engineering flow metrics, tax incentives and value-driven governance through OKRs and OrKRs. As an outcome, the work provides an implementation roadmap for innovation projects in Computer Science, focusing on hypothesis validation, scalability and economic-operational consistency.}
			
			%##################################################
			
			% Dados do guia
			%\begin{titlepage}
			%\pagestyle{empty}
			%\renewcommand{\baselinestretch}{1}
			%\enlargethispage{1.5cm}
			%\input{reitoria}
			%\cleardoublepage
			%\end{titlepage}
			
			%##################################################
			
			% descomente para habilitar a lista desejada
			\listoffigures                             % Lista de Figuras
			%\listofilustracoes
			%\listofgraficos							   % Lista de Gráficos
			\listoftables                              % Lista de Tabelas
			\listofquadros							   % Lista de Quadros
			%\listofexemplos
			%\listofteoremas
			\tableofcontents                           % Sumário
			
			\clearpage
			
			\pagestyle{unifal}
			
			%==============================================================================
			% incluindo os capitulos
			\chapter{INTRODUÇÃO}

O desenvolvimento de \textit{startups} de base tecnológica exige mais do que uma ideia inovadora: requer estruturação estratégica, coerência entre proposta de valor e mercado, e mecanismos de viabilidade econômica. Este trabalho sistematiza a proposta do Grupo 3 para a disciplina de empreendedorismo, com foco na construção de um modelo de negócio orientado à escalabilidade.

O problema central investigado é a transição entre a fase conceitual de uma solução tecnológica e sua consolidação como negócio viável. Para enfrentar esse desafio, foi adotada uma abordagem integrada baseada em cinco dimensões complementares: estratégica, estrutural, engenharia, fiscal e governança de valor.

Como objetivo geral, este documento apresenta uma arquitetura de modelagem capaz de apoiar decisões de produto, operações e crescimento da \textit{startup}. Como objetivos específicos, busca-se: (i) mapear elementos críticos para o desenho do \textit{Business Model Canvas}; (ii) integrar práticas modernas de gestão e engenharia; e (iii) descrever um roteiro de implementação com critérios de validação.

Além desta introdução, o Capítulo~\ref{cap:elementos} detalha o framework de cinco dimensões e seu diferencial competitivo. Por fim, o Capítulo~\ref{cap:conclusao} sintetiza os achados e apresenta os próximos passos para evolução do modelo.
			\chapter{DESENVOLVIMENTO DO MODELO DE NEGÓCIO}
\label{cap:elementos}

Este capítulo consolida o framework adotado pelo Grupo 3 para estruturar \textit{startups} de base tecnológica. A proposta integra cinco dimensões que, em conjunto, aumentam a probabilidade de construir um modelo de negócio robusto, escalável e financeiramente sustentável.

Em perspectiva executiva, o capítulo demonstra como decisões estratégicas, desenho organizacional e disciplina de execução convergem para transformar intenção empreendedora em resultado mensurável.

\section{Framework Integrado de Cinco Dimensões}

\subsection{Dimensão estratégica: PMBOK 7 e tailoring dinâmico}

Na dimensão estratégica, o foco está na geração de valor contínuo e no aprendizado validado. Em vez de processos rígidos, adota-se uma lógica orientada a princípios, com três eixos principais:

\begin{itemize}
    \item \textbf{Stewardship e governança participativa}: decisões orientadas ao valor entregue ao cliente e ao negócio;
    \item \textbf{Tailoring por maturidade}: adaptação dos ritos de gestão ao estágio da \textit{startup};
    \item \textbf{ROI mensurável}: cada ciclo de execução precisa produzir evidência de valor ou hipótese validada.
\end{itemize}

Como desdobramento operacional, essa dimensão se materializa em três frentes. Primeiro, define-se um ciclo curto de planejamento e revisão (quinzenal ou mensal), com hipóteses explícitas de mercado, produto e monetização. Em seguida, estabelece-se um quadro de priorização que classifica iniciativas por impacto esperado, esforço e risco de execução. Por fim, consolida-se uma rotina de \textit{review} estratégica, na qual os resultados do período são comparados com os objetivos definidos.

Para manter coerência entre estratégia e execução, recomenda-se monitorar: taxa de hipóteses validadas por ciclo, tempo médio entre decisão e experimento, e retorno incremental por iniciativa priorizada. Esses indicadores permitem ajustar rapidamente o direcionamento da \textit{startup} sem perder foco no valor entregue.

\subsection{Dimensão estrutural: platform design}

O modelo de negócio é concebido como um ecossistema multilateral, e não como cadeia linear de produção. Essa escolha amplia efeitos de rede e potencial de escala por meio de três motores:

\begin{itemize}
    \item \textbf{Motor de transações}: redução de fricção entre os participantes da plataforma;
    \item \textbf{Motor de aprendizado}: acumulação de dados e aumento de defensibilidade competitiva;
    \item \textbf{Feedback loops}: ciclos rápidos de retroalimentação entre produto, usuários e mercado.
\end{itemize}

Na prática, os desdobramentos estruturais envolvem mapear com precisão os atores do ecossistema (usuários finais, parceiros, fornecedores de tecnologia e canais de distribuição), definir a proposta de valor específica para cada lado da plataforma e estabelecer regras de interação que reduzam atrito de entrada. Esse desenho deve considerar mecanismos de atração inicial (\textit{seeding}), incentivos de permanência e barreiras de saída para aumentar retenção.

Também é necessário definir a lógica de monetização da plataforma desde cedo, mesmo que em versão preliminar. Entre os formatos mais usuais estão comissão por transação, assinaturas, \textit{tiers} de funcionalidades e serviços complementares. A escolha deve ser testada com indicadores como crescimento de base ativa, taxa de conversão e receita média por usuário.

\subsection{Dimensão de engenharia: Flow Framework}

A operação técnica é conectada diretamente ao desempenho de negócio. Para isso, o fluxo de trabalho é monitorado em quatro classes:

\begin{itemize}
    \item \textbf{Features}: entregas que geram valor percebido e potencial de receita;
    \item \textbf{Defects}: falhas que reduzem qualidade e experiência do usuário;
    \item \textbf{Riscos}: exposição a vulnerabilidades e não conformidades;
    \item \textbf{Dívida técnica}: parcela de capacidade reservada à sustentabilidade do código.
\end{itemize}

Recomenda-se reservar aproximadamente 20\% da capacidade da equipe para manutenção evolutiva e refatoração, mitigando crescimento desordenado da dívida técnica.

Como desdobramento de gestão técnica, cada classe de trabalho deve possuir política de entrada e critério de qualidade de saída. Em \textit{Features}, prioriza-se entrega incremental com validação de uso real; em \textit{Defects}, resposta proporcional ao impacto no cliente; em \textit{Riscos}, mitigação preventiva com monitoramento contínuo; e em dívida técnica, execução planejada para preservar velocidade futura de entrega.

Para conectar engenharia e negócio, recomenda-se acompanhar \textit{lead time}, \textit{throughput}, taxa de retrabalho e disponibilidade do serviço. A interpretação conjunta dessas métricas reduz decisões baseadas apenas em percepção e fortalece a previsibilidade da operação.

\subsection{Dimensão fiscal: Lei do Bem como alavanca de receita indireta}

A dimensão fiscal transforma atividades de pesquisa, desenvolvimento e inovação (PD\&I) em vantagem econômica por meio de conformidade estruturada. Com rastreabilidade adequada de iniciativas técnicas, torna-se possível recuperar parcela relevante dos custos de engenharia, com potencial de até 34\% conforme o enquadramento legal.

Os principais desdobramentos dessa dimensão são organizacionais e documentais. Do ponto de vista organizacional, é necessário delimitar quais frentes de trabalho se enquadram como PD\&I, com objetivos técnicos, incertezas e resultados esperados claramente registrados. Do ponto de vista documental, exige-se rastrear horas, artefatos de desenvolvimento, evidências de experimentação e resultados obtidos em cada ciclo.

Quando a governança fiscal é incorporada à rotina do time, o benefício não se limita ao incentivo tributário: há melhoria na disciplina de execução, maior clareza sobre custos de inovação e aumento da capacidade de justificar investimentos técnicos perante sócios e potenciais investidores.

\subsection{Governança de valor: VMO e OKRs}

A governança proposta substitui uma lógica centrada em controle burocrático por gestão orientada a valor. Assim, o \textit{Value Management Office} (VMO) organiza a priorização com uso de OKRs e OrKRs (\textit{Objectives, Risks and Key Results}), articulando estratégia, risco e mensuração contínua.

Como desdobramento prático, o VMO atua como camada de integração entre liderança, produto e engenharia. Sua função é tornar explícita a cadeia de valor de cada iniciativa: objetivo estratégico, hipótese de impacto, risco associado, métricas de acompanhamento e critério de continuidade ou interrupção.

Nessa lógica, os OKRs deixam de ser apenas metas declarativas e passam a orientar decisões de alocação de recursos. Já os OrKRs adicionam visibilidade sobre riscos críticos (tecnológicos, regulatórios, operacionais e de mercado), permitindo ação antecipada. O resultado esperado é maior disciplina de priorização e menor dispersão de esforço.

Como fechamento desta seção, o framework integrado posiciona o Grupo 3 para decidir com rapidez, aprender com evidências e preservar coerência entre estratégia, operação e desempenho econômico.

\section{Diferencial competitivo do Grupo 3}

O diferencial do Grupo 3 está na combinação de visão sistêmica e execução integrada. A proposta conecta decisões estratégicas, desenho de ecossistema, capacidade de engenharia e sustentabilidade financeira em uma arquitetura única de gestão.

\subsection{Diferencial pedagógico do modelo completo}

O modelo proposto fortalece a qualidade acadêmica do trabalho ao articular diferentes perspectivas de empreendedorismo em uma trilha única de decisão. Em vez de tratar os conteúdos de forma fragmentada, a abordagem integra formulação estratégica, validação prática e reflexão crítica sobre resultados.

Do ponto de vista formativo, a estrutura favorece o desenvolvimento simultâneo de competências analíticas, gerenciais e colaborativas. A autonomia discente é preservada na construção das hipóteses e na condução dos experimentos, enquanto a simulação de condições reais de mercado eleva a consistência das decisões apresentadas.

Outro ganho relevante é a combinação entre inovação incremental (aperfeiçoamento contínuo) e inovação de maior ruptura (novas propostas de valor), o que aumenta a robustez do raciocínio empreendedor e a aderência do projeto às demandas contemporâneas de negócio.

\subsection{Missão, entregáveis e foco do Grupo 3}

No arranjo geral do trabalho, o Grupo 3 tem como \textbf{missão} estruturar um negócio viável, convertendo ideias em um modelo com lógica econômica, coerência operacional e potencial de escala.

Para isso, os entregáveis são orientados por evidências e clareza de decisão: (i) \textit{Business Model Canvas} completo e internamente consistente; e (ii) modelagem objetiva de receitas, custos e parcerias-chave, com justificativas de priorização.

Esse foco amplia a qualidade técnica do capítulo ao explicitar a cadeia de valor do empreendimento, reforçando pensamento estratégico, estruturação organizacional e integração entre recursos, custos e monetização.

\subsection{Duplas funcionais e missões estratégicas}

Para viabilizar a execução do framework com alto grau de especialização e coordenação, o Grupo 3 organiza sua operação em três duplas complementares:

\begin{itemize}
    \item \textbf{Dupla A --- Arquitetos de Confiança (Hustler 1 + Hipster 1)}: tem como missão traduzir infraestrutura técnica em segurança como valor, criando um ecossistema no qual parceiros possam construir aplicações seguras sobre a base tecnológica proposta.
    \item \textbf{Dupla B --- Engenheiros de Defesa (Hacker 1 + Hipster 2)}: é responsável por produtizar scripts de defesa em módulos escaláveis de infraestrutura, assegurando também que a experiência de uso do \textit{dashboard} de segurança minimize fadiga de alertas e aumente a efetividade operacional.
    \item \textbf{Dupla C --- Guardiões do Caixa e Compliance (Hustler 2 + Hacker 2)}: atua na rastreabilidade das horas de PD\&I aplicadas a novos protocolos de segurança para viabilizar abatimento tributário de até 34\% via Lei do Bem, além de automatizar a geração de evidências técnicas exigidas pelo MCTI.
\end{itemize}

Como síntese gerencial, a Tabela~\ref{tab:trl-decisao-portfolio} consolida critérios de decisão por faixa de maturidade tecnológica, enquanto o Quadro~\ref{qua:governanca-integrada} articula a governança integrada entre tecnologia, finanças e conformidade.

Em termos de diferenciação competitiva, a proposta combina clareza estratégica, capacidade de coordenação e disciplina de governança, elevando a qualidade das decisões e a credibilidade do modelo perante avaliadores, parceiros e investidores.

\begin{table}[htb]
\centering
\caption{Matriz de maturidade e decisão para gestão do portfólio tecnológico.}
\label{tab:trl-decisao-portfolio}
\begin{tabular}{|p{0.15\textwidth}|p{0.35\textwidth}|p{0.38\textwidth}|}
\hline
\textbf{Faixa TRL} & \textbf{Foco principal} & \textbf{Decisão de gestão recomendada} \\
\hline
TRL 1--3 & Descoberta técnica e prova de conceito & Priorizar validação de hipótese crítica e reduzir escopo de incerteza. \\
\hline
TRL 4--6 & Prototipação e validação em ambiente relevante & Expandir testes, estruturar captação e preparar documentação de conformidade. \\
\hline
TRL 7--9 & Operação em ambiente real e escala & Acelerar comercialização, eficiência operacional e defesa competitiva. \\
\hline
\end{tabular}
\end{table}

\begin{table}[htb]
\centering
\caption{Quadro -- Governança para integração entre tecnologia, finanças e incentivos fiscais.}
\label{qua:governanca-integrada}
\begin{tabular}{|p{0.27\textwidth}|p{0.30\textwidth}|p{0.30\textwidth}|}
\hline
\textbf{Eixo de governança} & \textbf{Prática operacional} & \textbf{Valor agregado esperado} \\
\hline
Maturidade tecnológica & Revisão periódica de TRL/MRL com critério de avanço por evidência. & Menor risco de investimento prematuro e maior previsibilidade de entrega. \\
\hline
Engenharia financeira & Monitoramento de ROI/VPL por cenário e etapa de maturidade. & Melhor priorização de recursos e decisões de continuidade. \\
\hline
Incentivos fiscais & Rastreabilidade de horas, artefatos e resultados de PD\&I para Lei do Bem. & Economia tributária, segurança jurídica e reforço de caixa. \\
\hline
Conformidade regulatória & Registro de requisitos normativos e critérios de aceite por fase. & Redução de retrabalho e aceleração da entrada em mercado. \\
\hline
\end{tabular}
\end{table}

\section{Roteiro de implementação}

Nesta versão, o roteiro é tratado como \textbf{implementação executável} para uma startup de \textit{Hardware as a Service} (HaaS) em segurança da informação, com atuação inicial no Sul de Minas (Alfenas, Santa Rita do Sapucaí e Itajubá), priorizando agro e PMEs.

\subsection{Escopo implementado e tese de negócio}

O produto implementado é o \textbf{Kit Agro-Secure V1}: servidor de borda (\textit{home lab}), gateway de segurança, totem NFC e proteção elétrica, com monitoramento remoto de KPIs de continuidade operacional e segurança.

A tese de valor é objetiva: \textit{mesmo com internet instável e risco cibernético crescente, o cliente mantém operação crítica ativa, com rastreabilidade e governança de segurança em padrão corporativo}.

No posicionamento regional, a proposta aproveita a complementaridade do Sul de Minas: densidade do agronegócio (demanda operacional crítica), polos de engenharia eletrônica e telecomunicações (oferta técnica) e proximidade logística para atendimento presencial com SLA agressivo.

\subsection{Distribuição de responsabilidades entre as três duplas}

A incorporação integral do modelo Agro-Secure exige governança explícita de execução. Para reduzir sobreposição e aumentar accountability, as três duplas operam em frentes complementares, conforme a Tabela~\ref{tab:responsabilidades-duplas}.

\begin{table}[htb]
\centering
\caption{Distribuição de responsabilidades executivas entre as três duplas do Grupo 3.}
\label{tab:responsabilidades-duplas}
\begin{tabular}{|p{0.22\textwidth}|p{0.23\textwidth}|p{0.23\textwidth}|p{0.23\textwidth}|}
\hline
\textbf{Frente de execução} & \textbf{Dupla A (Arquitetos de Confiança)} & \textbf{Dupla B (Engenheiros de Defesa)} & \textbf{Dupla C (Caixa e Compliance)} \\
\hline
Arquitetura da oferta HaaS & Define padrão do kit, proposta técnica por segmento e desenho de valor percebido. & Especifica requisitos de firmware e operação segura do totem/edge. & Valida viabilidade econômica do padrão e impacto em CAPEX por cliente. \\
\hline
Instalação e operação em campo & Conduz diagnóstico inicial e desenho da topologia local. & Executa provisionamento, hardening e testes de failover no cliente. & Controla custo por instalação, deslocamento e produtividade por equipe. \\
\hline
SOC e monitoramento de KPIs & Define narrativa executiva dos relatórios para gestão do cliente. & Implementa coleta técnica (ameaças, uptime, temperatura, energia, link). & Consolida KPIs financeiros (MRR, payback, churn, margem de contribuição). \\
\hline
Validação comercial e escala & Estrutura pitch consultivo e parceiros regionais (cooperativas/associações). & Garante replicabilidade técnica dos kits e padrões de onboarding. & Modela cenários de expansão e necessidades de capital de giro. \\
\hline
Lei do Bem e conformidade & Enquadra hipóteses tecnológicas com foco em risco técnico real. & Produz evidências de experimentação, falhas e iterações técnicas. & Mantém timesheet, dossiê PD\&I e trilha para migração futura ao Lucro Real. \\
\hline
\end{tabular}
\end{table}

\subsection{Business Model Canvas completo}

\begin{table}[htb]
\centering
\caption{Business Model Canvas validável para startup InfoSec HaaS no Sul de Minas.}
\label{tab:bmc-haas-infosec}
\begin{tabular}{|p{0.30\textwidth}|p{0.62\textwidth}|}
\hline
\textbf{Bloco do BMC} & \textbf{Definição implementada} \\
\hline
Segmentos de clientes & Produtores rurais de café especial e grãos, cooperativas agrícolas e MMEs regionais do agro (laticínios, silos, comércio e logística). \\
\hline
Proposta de valor & Segurança \textit{enterprise} em formato \textit{plug and play}: servidores, roteadores seguros, totens NFC e monitoramento 24/7 sem necessidade de compra do hardware pelo cliente. \\
\hline
Canais & Venda consultiva in-loco, parcerias com cooperativas (ex.: Minasul e Cooxupé), feiras agropecuárias regionais e rede de indicação local. \\
\hline
Relacionamento com clientes & SLA agressivo, suporte presencial, onboarding operacional em campo e relatórios mensais executivos de ameaças bloqueadas, uptime preservado e estabilidade do ambiente. \\
\hline
Fontes de receita & Setup/instalação, mensalidade HaaS (hardware + firmware embarcado) e serviços de R\&D sob demanda para integrações e customizações críticas. \\
\hline
Recursos-chave & Estoque de hardware (servidor edge, gateway, NFC e energia), engenheiros de hardware/redes, SOC remoto em Alfenas e governança contábil para Lei do Bem. \\
\hline
Atividades-chave & Montagem e provisionamento de kits, instalação em campo (cabeamento/fixação), monitoramento de rede, manutenção e PD\&I em firmware e sincronização assíncrona segura. \\
\hline
Parcerias-chave & Fornecedores OEM (Brasil/Ásia), ISPs regionais, cooperativas agro, parceiros institucionais locais e escritório contábil com domínio em Lucro Real/Lei do Bem. \\
\hline
Estrutura de custos & CAPEX e depreciação de hardware, deslocamento/logística, folha técnica especializada, operação SOC 24/7, reposição de ativos e dispêndios de PD\&I. \\
\hline
\end{tabular}
\end{table}

\subsection{Funil de vendas e conversão regional}

Com base no anexo Agro-Secure, o funil comercial combina confiança local, prova técnica em campo e fechamento por evidência de valor econômico.

\begin{figure}[htb]
\centering
\fbox{\begin{minipage}{0.93\textwidth}
\textbf{Funil de vendas Agro-Secure (eixo Alfenas--Machado--Varginha)}\\
\textbf{Topo do funil (Dupla A):} prospecção por CNAE + introdução via cooperativas/sindicatos/associações.\\
\textbf{Meio do funil (Dupla B):} diagnóstico técnico e piloto de 30 dias com baseline de risco.\\
\textbf{Fundo do funil (Dupla A + Dupla C):} proposta HaaS com KPI executivo (uptime, ameaças bloqueadas, estabilidade) e contrato anual.\\
\textbf{Pós-venda e expansão (Dupla C + Dupla B):} renovação, upsell de integrações e escala por clusters regionais.
\end{minipage}}
\caption{Figura sugestiva do funil de vendas consultivo para conversão de pilotos em contratos HaaS.}
\label{fig:funil-vendas-haas}
\end{figure}

\begin{figure}[htb]
\centering
\fbox{\begin{minipage}{0.93\textwidth}
\textbf{Funil de metas de conversão (referência anual)}\\
\textbf{Etapa 1:} 15--20 prospects qualificados por território e CNAE.\\
\textbf{Etapa 2:} 3 pilotos ativos (30 dias) com coleta de KPIs operacionais.\\
\textbf{Etapa 3:} conversão comercial dos pilotos para contratos recorrentes.\\
\textbf{Etapa 4:} expansão via cooperativas (realista: 20 clientes no ano 1; otimista: 50).\\
\textbf{Meta financeira de referência:} MRR mês 12 entre R\$ 16.000,00 e R\$ 45.000,00.
\end{minipage}}
\caption{Figura sugestiva de funil de metas comerciais e financeiras para validação regional.}
\label{fig:funil-metas-conversao}
\end{figure}

\subsection{Arquitetura da solução e diagramas de validação}

As Figuras~\ref{fig:mapa-mental-haas} e~\ref{fig:usecase-campo} sintetizam, respectivamente, a arquitetura de valor e o fluxo operacional mínimo validado em campo.

\begin{figure}[htb]
\centering
\fbox{\begin{minipage}{0.92\textwidth}
\textbf{Mapa mental --- Oferta Agro-Secure V1}\\
\textbf{Núcleo:} Infraestrutura Segura HaaS\\
\textbf{Perímetro:} Gateway com firewall embarcado + controle de acesso NFC.\\
\textbf{Processamento local:} Servidor de borda para operação offline e sincronização assíncrona.\\
\textbf{Interface física:} Totem/Kiosk robusto para uso operacional em ambiente agro.\\
\textbf{Operação contínua:} SOC remoto com monitoramento de ameaças, energia, temperatura e disponibilidade.
\end{minipage}}
\caption{Figura sugestiva de mapa mental para arquitetura de oferta de valor.}
\label{fig:mapa-mental-haas}
\end{figure}

\begin{figure}[htb]
\centering
\fbox{\begin{minipage}{0.92\textwidth}
\textbf{Diagrama de caso de uso --- validação em campo}\\
Gestor da Fazenda $\rightarrow$ autentica via NFC no Totem Seguro.\\
Totem $\rightarrow$ valida acesso no Servidor Local HaaS.\\
Servidor Local $\rightarrow$ criptografa e sincroniza com nuvem/cooperativa.\\
SOC da startup $\dashrightarrow$ monitora anomalias, aplica atualização de firmware e registra KPIs.
\end{minipage}}
\caption{Figura sugestiva de caso de uso para operação local e monitoramento remoto.}
\label{fig:usecase-campo}
\end{figure}

\subsection{Backlog de validação e histórias de usuário}

O backlog prioriza \textbf{validação de disposição a pagar}, \textbf{adoção operacional} e \textbf{resiliência do kit} antes de qualquer escala de engenharia avançada.

\begin{table}[htb]
\centering
\caption{Backlog inicial orientado à validação de negócio (Epics e histórias).}
\label{tab:backlog-validacao}
\begin{tabular}{|p{0.20\textwidth}|p{0.20\textwidth}|p{0.52\textwidth}|}
\hline
\textbf{Épico} & \textbf{História} & \textbf{Critério de validação} \\
\hline
Disposição a pagar & HU 1.1 & Instalar 3 kits piloto em Alfenas por 30 dias e demonstrar ataques bloqueados, indisponibilidades evitadas e percepção de valor para conversão comercial. \\
\hline
Disposição a pagar & HU 1.2 & Entregar painel executivo mobile com métricas simples (ameaças bloqueadas, horas de operação preservadas e estabilidade da rede). \\
\hline
Adoção operacional & HU 2.1 & Prototipar totem com NFC e validar uso real por operadores de campo (luvas, poeira, rotina de turno), com taxa mínima de sucesso de autenticação. \\
\hline
Resiliência local & HU 3.1 & Garantir persistência local de eventos críticos (notas, pesagens, acessos) durante queda de link, com sincronização segura após retorno. \\
\hline
Conformidade PD\&I & HU 4.1 & Marcar tarefas inovadoras com etiqueta [PD\&I], registrar horas e evidências técnicas para formação contínua do dossiê Lei do Bem. \\
\hline
\end{tabular}
\end{table}

\subsection{Modelo técnico-operacional e unit economics do kit}

\begin{table}[htb]
\centering
\caption{Composição técnica e CAPEX médio do Kit Agro-Secure V1 por cliente.}
\label{tab:kit-capex-v1}
\begin{tabular}{|p{0.30\textwidth}|p{0.38\textwidth}|p{0.24\textwidth}|}
\hline
\textbf{Componente} & \textbf{Especificação de referência} & \textbf{Custo médio (R\$)} \\
\hline
Servidor Edge & Mini PC industrial fanless (16 GB RAM, SSD 500 GB) & 2.100,00 \\
\hline
Gateway de borda & Roteador corporativo (MikroTik/Ubiquiti) com VPN & 650,00 \\
\hline
Totem NFC & SBC + tela touch + leitor NFC + case reforçada & 1.500,00 \\
\hline
Energia & Nobreak senoidal puro 1200 VA & 1.200,00 \\
\hline
Miscelânea & Cabeamento, conectores, caixa de proteção e instalação & 450,00 \\
\hline
\textbf{CAPEX total por cliente} & \textbf{Kit físico alocado em comodato operacional} & \textbf{5.900,00} \\
\hline
\end{tabular}
\end{table}

\begin{table}[htb]
\centering
\caption{Precificação, payback e contribuição marginal do modelo HaaS.}
\label{tab:precificacao-payback}
\begin{tabular}{|p{0.46\textwidth}|p{0.34\textwidth}|}
\hline
\textbf{Indicador} & \textbf{Valor de referência} \\
\hline
Setup (instalação/auditoria) & R\$ 2.000,00 \\
\hline
Mensalidade HaaS & R\$ 850,00 \\
\hline
Receita no mês 1 por novo cliente & R\$ 2.850,00 \\
\hline
Exposição inicial de caixa por cliente & Aproximadamente R\$ 3.050,00 \\
\hline
Payback do hardware & Entre meses 5 e 6 \\
\hline
\end{tabular}
\end{table}

\subsection{Cenários de viabilidade e metas de validação}

\begin{table}[htb]
\centering
\caption{Cenários de viabilidade no primeiro ciclo anual de operação.}
\label{tab:cenarios-viabilidade}
\begin{tabular}{|p{0.25\textwidth}|p{0.20\textwidth}|p{0.20\textwidth}|p{0.20\textwidth}|}
\hline
\textbf{Métrica} & \textbf{Pessimista} & \textbf{Realista} & \textbf{Otimista} \\
\hline
Clientes ativos no ano 1 & 5 & 20 & 50 \\
\hline
CAPEX acumulado (R\$) & 20.000 & 80.000 & 180.000 \\
\hline
MRR no mês 12 (R\$) & 4.000 & 16.000 & 45.000 \\
\hline
Leitura executiva & Pivotar escopo para preservar caixa & Crescimento orgânico com estabilização & Escalar com capital externo \\
\hline
\end{tabular}
\end{table}

\subsection{Tratamento estatístico bayesiano da viabilidade}

Para integrar incerteza de mercado, execução e risco de capital intensivo, adota-se um modelo bayesiano de atualização de probabilidade de sucesso do negócio. Considera-se sucesso como atingir tração comercial com manutenção de continuidade operacional e evolução para equilíbrio econômico no horizonte de 12 a 14 meses.

\begin{table}[htb]
\centering
\caption{Evidências usadas na atualização bayesiana da probabilidade de sucesso.}
\label{tab:bayes-evidencias}
\begin{tabular}{|p{0.29\textwidth}|p{0.20\textwidth}|p{0.20\textwidth}|p{0.21\textwidth}|}
\hline
\textbf{Evidência observada} & \textbf{P(E$\mid$Sucesso)} & \textbf{P(E$\mid$Fracasso)} & \textbf{Razão de verossimilhança} \\
\hline
Demanda agro regional aquecida & 0,80 & 0,40 & 2,00 \\
\hline
Canal cooperativista estruturado & 0,75 & 0,45 & 1,67 \\
\hline
Dor técnica elevada (conectividade/ataques) & 0,85 & 0,50 & 1,70 \\
\hline
Restrição de crédito rural (evidência adversa) & 0,40 & 0,70 & 0,57 \\
\hline
Capacidade interna de execução (3 duplas) & 0,70 & 0,50 & 1,40 \\
\hline
\end{tabular}
\end{table}

Com probabilidade \textit{a priori} de 50\%, a atualização com as evidências da Tabela~\ref{tab:bayes-evidencias} resulta em probabilidade \textit{a posteriori} aproximada de \textbf{81,9\%} para sucesso do modelo, mantendo disciplina de risco operacional e financeiro.

\begin{table}[htb]
\centering
\caption{Atualização bayesiana dos cenários e MRR esperado no mês 12.}
\label{tab:bayes-cenarios-mrr}
\begin{tabular}{|p{0.20\textwidth}|p{0.16\textwidth}|p{0.18\textwidth}|p{0.18\textwidth}|p{0.16\textwidth}|}
\hline
\textbf{Cenário} & \textbf{Priors} & \textbf{P(E$\mid$Cenário)} & \textbf{Posterior} & \textbf{MRR mês 12 (R\$)} \\
\hline
Pessimista & 0,30 & 0,40 & 0,203 & 4.000 \\
\hline
Realista & 0,50 & 0,70 & 0,593 & 16.000 \\
\hline
Otimista & 0,20 & 0,60 & 0,203 & 45.000 \\
\hline
\multicolumn{4}{|r|}{\textbf{MRR esperado (valor esperado bayesiano)}} & \textbf{19.458} \\
\hline
\end{tabular}
\end{table}

Esse tratamento estatístico reforça a estratégia de execução por ondas: preserva prudência de caixa no curto prazo e, ao mesmo tempo, sustenta decisão de escala progressiva apoiada em evidências observáveis de conversão e retenção.

Como métrica de execução incremental, adota-se uma trilha em três ondas: (i) pilotos com 3 kits e conversão por valor demonstrado; (ii) expansão via parcerias com cooperativas; e (iii) consolidação regional com operação SOC 24/7 e backlog técnico escalável.

\begin{quadro}[htb]
\centering
\caption{Plano de execução por ondas com protagonismo das duplas.}
\label{qua:ondas-execucao-duplas}
\begin{tabular}{|p{0.20\textwidth}|p{0.20\textwidth}|p{0.16\textwidth}|p{0.14\textwidth}|p{0.20\textwidth}|}
\hline
\textbf{Onda} & \textbf{Meta de negócio} & \textbf{Dupla líder} & \textbf{Prazo de referência} & \textbf{Entregável crítico} \\
\hline
Onda 1 -- Pilotos & Validar conversão e disposição a pagar em campo & Dupla A & Meses 1--4 & 3 pilotos ativos, relatório executivo de valor e taxa de conversão monitorada \\
\hline
Onda 2 -- Escala regional & Acelerar aquisição com cooperativas e padronizar implantação & Dupla B & Meses 5--8 & Playbook técnico de instalação, onboarding e operação resiliente \\
\hline
Onda 3 -- Consolidação & Elevar MRR, governança e prontidão para expansão/captação & Dupla C & Meses 9--12 & Painel financeiro-operacional e dossiê PD\&I auditável \\
\hline
\end{tabular}
\end{quadro}

\begin{quadro}[htb]
\centering
\caption{Hipóteses críticas e indicadores de validação de negócio.}
\label{qua:hipoteses-validacao}
\begin{tabular}{|p{0.30\textwidth}|p{0.26\textwidth}|p{0.30\textwidth}|}
\hline
\textbf{Hipótese} & \textbf{Indicador-chave} & \textbf{Critério de sucesso} \\
\hline
Cliente paga pela continuidade operacional & Taxa de conversão piloto$\rightarrow$contrato & Conversão mínima de 40\% em até 45 dias \\
\hline
Kit reduz risco operacional real & Horas de operação preservadas em quedas de link & Evidência mensal de indisponibilidade evitada \\
\hline
Valor de segurança é compreendido pela gestão & Relatório de ameaças bloqueadas aceito na renovação & Renovação no mês 2 superior a 80\% \\
\hline
Modelo sustenta expansão regional & Relação MRR/CAPEX por cluster geográfico & Ponto de equilíbrio operacional até mês 14 \\
\hline
\end{tabular}
\end{quadro}

\subsection{Lei do Bem como eixo de vantagem competitiva}

A Lei do Bem (Lei n\textordmasculine\ 11.196/2005) é tratada como mecanismo estratégico de médio prazo. O desenho operacional separa atividades rotineiras de instalação das atividades de \textbf{risco tecnológico}, que são as elegíveis para incentivo fiscal quando houver migração sustentável ao regime de Lucro Real.

\begin{table}[htb]
\centering
\caption{Estrutura mínima do dossiê técnico de PD\&I para Lei do Bem.}
\label{tab:dossie-lei-bem}
\begin{tabular}{|p{0.32\textwidth}|p{0.58\textwidth}|}
\hline
\textbf{Seção do dossiê} & \textbf{Implementação prática na startup} \\
\hline
Escopo do projeto & Definição de objetivo técnico, período, incerteza e critério de sucesso por projeto embarcado. \\
\hline
Risco tecnológico & Registro do estado da arte, barreira técnica e hipótese de superação em ambiente agro de alta latência. \\
\hline
Evidência experimental & Prototipagem, teste de campo, falhas registradas, iterações e decisão técnica por versão. \\
\hline
Timesheet técnico & Apontamento de horas por tarefa [PD\&I], separando engenharia inovadora de operação rotineira. \\
\hline
Dispêndios elegíveis & Consolidação de mão de obra técnica, materiais de protótipo e custos vinculados a experimentação. \\
\hline
\end{tabular}
\end{table}

Como fechamento executivo desta seção, a implementação proposta transforma o roteiro em \textbf{plano de execução validável}, com entrega concreta de valor ao cliente, disciplina financeira para um modelo intensivo em hardware e construção antecipada de conformidade para capturar benefícios da Lei do Bem no momento econômico adequado.

Em síntese estratégica, a integração entre as três duplas permite converter complexidade técnica em resultado de negócio: Dupla A acelera tração e posicionamento, Dupla B assegura excelência operacional em campo e Dupla C protege sustentabilidade financeira e vantagem fiscal de longo prazo.

\section{Avaliação}

Para tornar os critérios de desempenho transparentes e comparáveis entre equipes, adota-se a Tabela~\ref{tab:criterios-avaliacao}, com pontuação total de 10,0 pontos.

Para maximizar desempenho em todos os critérios, recomenda-se que cada seção entregue três evidências mínimas: problema de negócio claramente delimitado, decisão metodológica justificada e resultado verificável (métrica, validação ou aprendizado documentado). Essa lógica reduz subjetividade na correção e fortalece a consistência da apresentação final.

Como fechamento avaliativo, a estratégia recomendada é sustentar cada afirmação com evidência objetiva, explicitando nexo entre problema, método e resultado, de modo a maximizar clareza, rigor e impacto da apresentação.

\begin{table}[htb]
\centering
\caption{Matriz de atendimento dos critérios de avaliação com argumentação de contrapartida.}
\label{tab:criterios-avaliacao}
\begin{tabular}{|p{0.23\textwidth}|p{0.39\textwidth}|p{0.24\textwidth}|p{0.08\textwidth}|}
\hline
\textbf{Critério} & \textbf{Argumento de contrapartida (como o trabalho atende)} & \textbf{Evidência objetiva no documento} & \textbf{Pontos} \\
\hline
Clareza da proposta & O problema regional, o ICP por CNAE e a tese de valor HaaS foram definidos com recorte geográfico e econômico explícito. & BMC validável, matriz de prospecção, funil com etapas e metas de conversão. & 1,0 \\
\hline
Aplicação dos conceitos & Os referenciais da disciplina foram operacionalizados em artefatos de gestão, engenharia, finanças e governança de valor. & Framework de cinco dimensões, backlog, hipóteses, cenários e plano por ondas. & 1,5 \\
\hline
Inovação e criatividade & A proposta integra hardware, firmware embarcado, edge computing e Lei do Bem como diferencial técnico-fiscal. & Kit Agro-Secure V1, arquitetura de operação offline, trilha de PD\&I e dossiê técnico. & 2,0 \\
\hline
Qualidade dos entregáveis & O material apresenta coerência executiva entre estratégia, operação, modelagem financeira e conformidade regulatória. & Tabelas de CAPEX/payback, cenários de viabilidade, quadros de responsabilidades e governança. & 1,5 \\
\hline
Apresentação e participação & A narrativa está estruturada para defesa oral objetiva, com argumentos verificáveis e divisão clara de protagonismo entre duplas. & Matriz de responsabilidades, funil comercial revisado, indicadores de validação e sínteses executivas por seção. & 4,0 \\
\hline
\textbf{Total} &  &  & \textbf{10,0} \\
\hline
\end{tabular}
\end{table}
			\chapter{CONCLUSÃO}
\label{cap:conclusao}

Este trabalho apresentou a estruturação do modelo de negócio do Grupo 3 com base em uma abordagem integrada para \textit{startups} de base tecnológica. A análise evidencia que a combinação entre estratégia adaptativa, design de plataforma, engenharia orientada a fluxo, inteligência fiscal e governança por valor fortalece a viabilidade do empreendimento desde as fases iniciais.

Como principal contribuição, o estudo organiza um caminho prático para transformar ideias em modelos de negócio testáveis e escaláveis, reduzindo incertezas típicas do processo empreendedor. Destaca-se, ainda, a importância de conectar métricas técnicas e métricas econômicas para orientar decisões de priorização com maior consistência.

Como trabalhos futuros, recomenda-se aplicar o framework a um problema real de mercado, construir o canvas completo da solução, realizar modelagem financeira com cenários e conduzir ciclos de validação com usuários e especialistas. Esses passos permitem consolidar o modelo proposto e avaliar seu potencial de implementação em contexto real.

			
			%==============================================================================
			% Incluindo bibliografia
			%\bibliographystyle{plain}             % estilo para labels em numeros
			%\bibliographystyle{alpha}             % estilo para labels em iniciais
			\bibliographystyle{abntex2-alf}           % estilo para referências usando ABNT, 
			% precisa instalar o abntex para usar!!!
			%\renewcommand*{\mkbibnamefamily}[1]{#1}%                                       
			
			%inclui Referências Bibliográficas
			\referencias
			\nocite{wadhwani2026startups,sebrae2026modelos,sebrae2026conheca,sebrae2026modeloavancado,oconnor2018limits,sinova2026leidobem,abgi2026leiinformatica,taxgroup2026leidobem,inpi2026trl,finep2024maisinovacao,finep2022trl,embrapa2026trl,finep2026anexo5,expertzy2026finep,abgi2026trlleidobem,finep2023avaliacao,brandao2023trl,gao2018timing,weshine2026roi,cubo2026roi,venturehub2026roi,mercadopago2026valuation,ey2026valuation,qubit2026dcf,phoenix2026dcf,valutico2026valuation,li2024valuation,macke2026indicadores,jotage2026contabilidade,abgi2026leidobembeneficios,bellver2026planejamento,abgi2026startupsleidobem,systemiq2025diversity,portaldaindustria2026leidobem,deavila2018leidobem,ipea2020efetividade,ipea2020efetividadehandle,mcti2020guia,mcti2017guia,tigre2017gestao}
			%\printbibliography%
			\bibliography{refbib}			% arquivo exemplo refbib.bib
			%==============================================================================
			% Incluindo anexos num1erados com letras maiusculas.
			%\apendices
			\apendice{Esquemas visuais do modelo do Grupo 3}
\label{cap:apendice}

Este apêndice reúne diagramas vetoriais do framework, com foco no objetivo primordial do Grupo 3: segurança como valor, escalabilidade técnica, compliance e retorno econômico.

\renewcommand{\thesection}{\arabic{section}}

\section{Diagrama integrado das cinco dimensões}
\begin{figure}[htb]
\centering
\begin{tikzpicture}[
node distance=1.4cm and 1.8cm,
box/.style={draw, rounded corners=4pt, thick, align=center, minimum width=3.0cm, minimum height=1.1cm, fill=gray!12},
arr/.style={-Latex, thick}
]
\node[box] (a) {Estratégica\\PMBOK 7 + Tailoring};
\node[box, right=of a] (b) {Estrutural\\Platform Design};
\node[box, right=of b] (c) {Engenharia\\Flow Framework};
\node[box, below=of b] (d) {Fiscal\\Lei do Bem};
\node[box, below=of c] (e) {Governança\\VMO + OKRs/OrKRs};
\draw[arr] (a)--(b);
\draw[arr] (b)--(c);
\draw[arr] (c)--(e);
\draw[arr] (d)--(e);
\draw[arr] (a)|-(d);
\draw[arr] (e)-|(a);
\end{tikzpicture}
\caption{Integração sistêmica das cinco dimensões.}
\end{figure}

\section{Mapa mental estratégico revisado}
\begin{figure}[htb]
\centering
\begin{tikzpicture}[
node distance=1.4cm and 1.7cm,
box/.style={draw, rounded corners=4pt, thick, align=center, minimum width=3.3cm, minimum height=1.0cm, fill=gray!12},
arr/.style={-Latex, thick}
]
\node[box, minimum width=4.0cm] (core) {Agro-Secure HaaS\\núcleo de valor};
\node[box, above left=of core] (m1) {Mercado-alvo\\Agro + PMEs};
\node[box, above=of core] (m2) {Oferta\\segurança enterprise};
\node[box, above right=of core] (m3) {Arquitetura\\edge + NFC + SOC};
\node[box, below left=of core] (m4) {Economia\\CAPEX 5,9k / payback 5};
\node[box, below=of core] (m5) {Conformidade\\PD\&I + Lei do Bem};
\node[box, below right=of core] (m6) {3 duplas\\A tração, B operação, C caixa};
\draw[arr] (core)--(m1);
\draw[arr] (core)--(m2);
\draw[arr] (core)--(m3);
\draw[arr] (core)--(m4);
\draw[arr] (core)--(m5);
\draw[arr] (core)--(m6);
\end{tikzpicture}
\caption{Mapa mental estratégico com foco em validação regional, execução e conformidade.}
\end{figure}

\section{Item 2 do anexo --- Business Model Canvas}
\noindent Este bloco consolida visualmente os nove componentes do \textit{Business Model Canvas} aderente ao anexo Agro-Secure.
\begin{figure}[H]
\centering
\begin{tikzpicture}[x=0.98cm,y=0.98cm]
\draw[thick] (0,0) rectangle (16,9.6);
\draw[thick] (3,0)--(3,9.6);
\draw[thick] (7,0)--(7,9.6);
\draw[thick] (11,0)--(11,9.6);
\draw[thick] (13.5,0)--(13.5,9.6);
\draw[thick] (0,4.8)--(13.5,4.8);
\draw[thick] (3,2.4)--(7,2.4);
\draw[thick] (11,2.4)--(13.5,2.4);

\node[anchor=north west,font=\bfseries\scriptsize] at (0.1,9.5) {Parcerias-chave};
\node[anchor=north west,font=\scriptsize,align=left,text width=2.8cm] at (0.1,9.0) {\textbullet{} OEMs (Brasil/Ásia)\\\textbullet{} ISPs regionais\\\textbullet{} Cooperativas\\\textbullet{} Contabilidade (Lucro Real)};

\node[anchor=north west,font=\bfseries\scriptsize] at (3.1,9.5) {Atividades-chave};
\node[anchor=north west,font=\scriptsize,align=left,text width=3.8cm] at (3.1,9.0) {\textbullet{} Montagem e provisionamento\\\textbullet{} Instalação em campo\\\textbullet{} Monitoramento SOC\\\textbullet{} R\&D de firmware};

\node[anchor=north west,font=\bfseries\scriptsize] at (3.1,4.5) {Recursos-chave};
\node[anchor=north west,font=\scriptsize,align=left,text width=3.8cm] at (3.1,4.0) {\textbullet{} Kit Agro-Secure V1\\\textbullet{} 3 duplas funcionais\\\textbullet{} SOC em Alfenas};

\node[anchor=north west,font=\bfseries\scriptsize] at (7.1,9.5) {Proposta de valor};
\node[anchor=north west,font=\scriptsize,align=left,text width=3.8cm] at (7.1,9.0) {\textbullet{} Segurança enterprise plug and play\\\textbullet{} Operação contínua mesmo offline\\\textbullet{} Sem compra de hardware};

\node[anchor=north west,font=\bfseries\scriptsize] at (11.1,9.5) {Relacionamento};
\node[anchor=north west,font=\scriptsize,align=left,text width=2.3cm] at (11.1,9.0) {\textbullet{} SLA agressivo\\\textbullet{} Suporte presencial\\\textbullet{} Relatório mensal};

\node[anchor=north west,font=\bfseries\scriptsize] at (11.1,4.5) {Canais};
\node[anchor=north west,font=\scriptsize,align=left,text width=2.3cm] at (11.1,4.0) {\textbullet{} Venda consultiva\\\textbullet{} Feiras agro\\\textbullet{} Indicação local};

\node[anchor=north west,font=\bfseries\scriptsize] at (13.6,9.5) {Segmentos};
\node[anchor=north west,font=\scriptsize,align=left,text width=2.3cm] at (13.6,9.0) {\textbullet{} Produtores\\\textbullet{} Cooperativas\\\textbullet{} MMEs agro};

\node[anchor=north west,font=\bfseries\scriptsize] at (0.1,2.1) {Estrutura de custos};
\node[anchor=north west,font=\scriptsize,align=left,text width=7.6cm] at (0.1,1.6) {\textbullet{} CAPEX/depreciação do hardware\quad \textbullet{} Logística e deslocamento\quad \textbullet{} Folha técnica e SOC\quad \textbullet{} PD\&I};

\node[anchor=north west,font=\bfseries\scriptsize] at (8.1,2.1) {Fontes de receita};
\node[anchor=north west,font=\scriptsize,align=left,text width=5.2cm] at (8.1,1.6) {\textbullet{} Setup (instalação)\quad \textbullet{} Mensalidade HaaS (hardware + firmware)\quad \textbullet{} Serviços de R\&D};
\end{tikzpicture}
\caption{BMC tradicional preenchido com base no anexo Agro-Secure.}
\end{figure}

\section{Item 3 do anexo --- Economia unitária e retorno}
\noindent A economia unitária do modelo HaaS evidencia o equilíbrio entre investimento inicial em hardware e recorrência contratual.
\begin{table}[H]
\centering
\caption{Economia unitária do Kit Agro-Secure V1.}
\begin{tabular}{|p{6.0cm}|p{5.5cm}|}
\hline
\textbf{Indicador} & \textbf{Referência de validação} \\
\hline
CAPEX por kit & R\$ 5.900,00 \\
\hline
Mensalidade recorrente (HaaS) & R\$ 850,00 por cliente/mês \\
\hline
Setup de instalação & R\$ 2.000,00 \\
\hline
Payback do hardware & Aproximadamente mês 5 \\
\hline
Leitura de margem (mês 6 ao 12) & Mensalidade com alta contribuição marginal \\
\hline
\end{tabular}
\end{table}

\section{Item 4 do anexo --- Arquitetura de infraestrutura segura}
\noindent A arquitetura prioriza resiliência local no cliente e monitoramento contínuo remoto via SOC.
\begin{figure}[H]
\centering
\begin{tikzpicture}[
node distance=1.2cm and 1.8cm,
box/.style={draw, rounded corners=4pt, thick, align=center, minimum width=4.0cm, minimum height=1.0cm, fill=gray!12},
arr/.style={-Latex, thick}
]
\node[box] (p) {Perímetro seguro\\Gateway + Firewall + NFC};
\node[box, right=of p] (l) {Processamento local\\Servidor Edge (Home Lab)};
\node[box, right=of l] (i) {Interface física\\Totem/Kiosk};
\node[box, below=of l] (s) {SOC remoto (Alfenas)\\Monitoramento 24/7};
\draw[arr] (p)--(l);
\draw[arr] (l)--(i);
\draw[arr] (s)--(l);
\draw[arr] (s)-| (i);
\end{tikzpicture}
\caption{Arquitetura operacional da oferta HaaS de segurança da informação.}
\end{figure}

\section{Figura 5 (principal) --- Funil de vendas revisado}
\noindent A Figura 5 representa o funil operacional completo, da prospecção por CNAE até a escala regional por cooperativas.
\begin{figure}[H]
\centering
\begin{tikzpicture}
\node[draw,trapezium,trapezium left angle=80,trapezium right angle=100,minimum width=10.6cm,minimum height=1cm,fill=gray!12] (f1) at (0,0) {Prospects qualificados por CNAE (15--20)};
\node[draw,trapezium,trapezium left angle=80,trapezium right angle=100,minimum width=8.9cm,minimum height=1cm,fill=gray!12] (f2) at (0,-1.25) {Diagnóstico técnico e auditoria de risco (7 dias)};
\node[draw,trapezium,trapezium left angle=80,trapezium right angle=100,minimum width=7.2cm,minimum height=1cm,fill=gray!12] (f3) at (0,-2.5) {Pilotos ativos (3 kits / 30 dias)};
\node[draw,trapezium,trapezium left angle=80,trapezium right angle=100,minimum width=5.6cm,minimum height=1cm,fill=gray!12] (f4) at (0,-3.75) {Conversão para HaaS recorrente};
\node[draw,trapezium,trapezium left angle=80,trapezium right angle=100,minimum width=4.0cm,minimum height=1cm,fill=gray!12] (f5) at (0,-5.0) {Escala via cooperativas};
\draw[-Latex,thick] (f1.south)--(f2.north);
\draw[-Latex,thick] (f2.south)--(f3.north);
\draw[-Latex,thick] (f3.south)--(f4.north);
\draw[-Latex,thick] (f4.south)--(f5.north);
\node[draw,rounded corners=2pt,align=left,font=\scriptsize,fill=gray!8,text width=10.9cm] at (0,-6.45) {
\textbf{Papéis por etapa:} Dupla A (segmentação e abordagem consultiva); Dupla B (baseline técnico, implantação e coleta de KPIs); Dupla A+Dupla C (proposta, contrato e viabilidade); Dupla C (governança financeira e escala regional).};
\end{tikzpicture}
\caption{Funil de vendas orientado por validação de valor e protagonismo das três duplas.}
\end{figure}

\section{Item 6 do anexo --- Lei do Bem como barreira de entrada}
\noindent A estratégia fiscal é tratada como vantagem competitiva progressiva, com disciplina de evidências desde o início da operação.
\begin{table}[H]
\centering
\caption{Trilha de conformidade para captura futura dos incentivos da Lei do Bem.}
\begin{tabular}{|p{4.5cm}|p{7.0cm}|}
\hline
\textbf{Eixo} & \textbf{Implementação prática} \\
\hline
Documentação rigorosa & Timesheets, backlog [PD\&I], registro de falhas técnicas e relatórios de incerteza desde o dia 1. \\
\hline
Separação do inovador vs rotineiro & Distinção entre instalação operacional e desenvolvimento de firmware/orquestração offline. \\
\hline
Preparação para Lucro Real & Estruturação contábil e técnica para migração no momento de maturidade econômica (meta: ano 3). \\
\hline
Otimização fiscal esperada & Potencial de abatimento entre 60\% e 80\% sobre dispêndios elegíveis de PD\&I. \\
\hline
\end{tabular}
\end{table}

\section{Item 7 do anexo --- Arquitetura e componentes do kit}
\noindent A tabela a seguir explicita os componentes físicos mínimos para execução do modelo em campo com robustez operacional.
\begin{table}[H]
\centering
\caption{Composição técnica do Kit Agro-Secure V1 e função operacional.}
\begin{tabular}{|p{3.3cm}|p{5.1cm}|p{3.2cm}|}
\hline
\textbf{Componente} & \textbf{Função na operação} & \textbf{Custo médio} \\
\hline
Servidor Edge & Processamento local, persistência offline e sincronização segura de dados críticos. & R\$ 2.100,00 \\
\hline
Gateway de borda & Firewall primário, VPN e proteção do perímetro de rede local. & R\$ 650,00 \\
\hline
Totem NFC & Controle de acesso operacional sem uso de credenciais frágeis ou pendrives. & R\$ 1.500,00 \\
\hline
Nobreak senoidal & Proteção contra oscilações elétricas e continuidade da infraestrutura. & R\$ 1.200,00 \\
\hline
Miscelânea de instalação & Cabeamento, conectores e elementos de fixação/proteção do ambiente. & R\$ 450,00 \\
\hline
\textbf{CAPEX total por kit} & \textbf{Infraestrutura completa em comodato HaaS} & \textbf{R\$ 5.900,00} \\
\hline
\end{tabular}
\end{table}

\section{Painel estatístico do agronegócio mineiro (2024--2026)}
\noindent Esta seção consolida os dados macroeconômicos, cooperativistas e de risco cibernético utilizados para sustentar a viabilidade do modelo HaaS no Sul de Minas.

\begin{table}[H]
\centering
\caption{Indicadores macroeconômicos de Minas Gerais com impacto na tese Agro-Secure.}
\begin{tabular}{|p{6.4cm}|p{5.1cm}|}
\hline
\textbf{Indicador (2025)} & \textbf{Valor} \\
\hline
Receita de exportação do estado & US\$ 18,1 bilhões \\
\hline
Participação do agro nas exportações mineiras & 43,7\% \\
\hline
VBP agropecuário total de Minas Gerais & R\$ 171,80 bilhões \\
\hline
VBP do café em Minas Gerais & R\$ 57,52 bilhões \\
\hline
VBP do leite em Minas Gerais & R\$ 26,73 bilhões \\
\hline
Participação de Minas no café nacional & 54\% \\
\hline
\end{tabular}
\end{table}

\begin{table}[H]
\centering
\caption{Base cooperativista estratégica no Sul de Minas.}
\begin{tabular}{|p{2.6cm}|p{3.3cm}|p{2.9cm}|p{3.0cm}|}
\hline
\textbf{Cooperativa} & \textbf{Faturamento} & \textbf{Volume de café} & \textbf{Base de cooperados} \\
\hline
Cooxupé & R\$ 10,7 bilhões & 6,1 milhões de sacas & \textasciitilde 18.000 famílias \\
\hline
Minasul & Consolidado 2025 & 1,8 milhão de sacas (recebido) & 9.000+ cooperados \\
\hline
Cocatrel & R\$ 3,45 bilhões & 1,83 milhão de sacas & Base Sul de Minas \\
\hline
\end{tabular}
\end{table}

\begin{table}[H]
\centering
\caption{Indicadores de infraestrutura e cibersegurança com efeito direto na viabilidade.}
\begin{tabular}{|p{6.4cm}|p{5.1cm}|}
\hline
\textbf{Variável de risco/operação} & \textbf{Valor de referência} \\
\hline
Produtores com internet de qualidade (MG) & 17\% \\
\hline
Ataques cibernéticos no agro brasileiro (2025) & 39.034 \\
\hline
Média mensal de tentativas de ataque (2025) & 3,2 mil/mês \\
\hline
Participação estimada de ransomware nos incidentes & $>$50\% \\
\hline
Retração estimada do crédito rural em 2025 & -16\% \\
\hline
\end{tabular}
\end{table}

\section{Gráficos estratégicos (pizza, radar e barras)}

\begin{figure}[H]
\centering
\begin{tikzpicture}[scale=1.0]
\def\r{2.3}
\fill[gray!55] (0,0) -- (0:\r) arc (0:202:\r) -- cycle;
\fill[gray!35] (0,0) -- (202:\r) arc (202:258:\r) -- cycle;
\fill[gray!20] (0,0) -- (258:\r) arc (258:295:\r) -- cycle;
\fill[gray!75] (0,0) -- (295:\r) arc (295:360:\r) -- cycle;
\draw[thick] (0,0) circle (\r);
\node[font=\scriptsize,align=left] at (4.0,1.9) {\textbf{Legenda (exportações agro MG)}\\Café: 56,1\%\\Soja: 15,6\%\\Sucroalcooleiro: 10,3\%\\Outros: 18,0\%};
\end{tikzpicture}
\caption{Gráfico de pizza da composição percentual das exportações agro de Minas Gerais.}
\end{figure}

\begin{figure}[H]
\centering
\begin{tikzpicture}
\begin{axis}[
    ybar,
    width=0.86\textwidth,
    height=6.1cm,
    ymin=0,
    ymax=65,
    ylabel={R\$ bilhões},
    symbolic x coords={Café,Leite,Pecuária},
    xtick=data,
    nodes near coords,
    nodes near coords align={vertical},
    bar width=20pt,
    grid=major
]
\addplot[fill=gray!35] coordinates {(Café,57.52) (Leite,26.73) (Pecuária,17.98)};
\end{axis}
\end{tikzpicture}
\caption{Gráfico de barras do VBP mineiro por cadeias relevantes para a tese Agro-Secure.}
\end{figure}

\begin{figure}[H]
\centering
\begin{tikzpicture}[scale=1.05]
\def\R{3.0}
% grade hexagonal
\foreach \k in {0.2,0.4,0.6,0.8,1.0}{
  \draw[gray!35] ({\k*\R*cos(90)},{\k*\R*sin(90)}) -- ({\k*\R*cos(30)},{\k*\R*sin(30)}) -- ({\k*\R*cos(-30)},{\k*\R*sin(-30)}) -- ({\k*\R*cos(-90)},{\k*\R*sin(-90)}) -- ({\k*\R*cos(-150)},{\k*\R*sin(-150)}) -- ({\k*\R*cos(150)},{\k*\R*sin(150)}) -- cycle;
}
% eixos
\foreach \a in {90,30,-30,-90,-150,150}{
  \draw[gray!55] (0,0) -- ({\R*cos(\a)},{\R*sin(\a)});
}
% polígono dos dados (90,85,80,45,70,78)
\draw[thick,fill=gray!35,fill opacity=0.35]
({0.90*\R*cos(90)},{0.90*\R*sin(90)}) --
({0.85*\R*cos(30)},{0.85*\R*sin(30)}) --
({0.80*\R*cos(-30)},{0.80*\R*sin(-30)}) --
({0.45*\R*cos(-90)},{0.45*\R*sin(-90)}) --
({0.70*\R*cos(-150)},{0.70*\R*sin(-150)}) --
({0.78*\R*cos(150)},{0.78*\R*sin(150)}) -- cycle;
% rótulos
\node[font=\scriptsize,align=center] at (0,3.35) {Demanda};
\node[font=\scriptsize,align=left] at (3.0,1.8) {Canal coop.};
\node[font=\scriptsize,align=left] at (3.0,-1.8) {Risco ciber};
\node[font=\scriptsize,align=center] at (0,-3.35) {Conectividade};
\node[font=\scriptsize,align=right] at (-3.0,-1.8) {Prontidão fiscal};
\node[font=\scriptsize,align=right] at (-3.0,1.8) {Execução técnica};
\end{tikzpicture}
\caption{Radar de prontidão estratégica do modelo HaaS no Sul de Minas (escala 0--100).}
\end{figure}

\section{Curvas de ROI acumulado da validação}
\begin{equation}
\mathrm{ROI}(\%) = \frac{\text{Benefício anual esperado} - \text{Investimento total}}{\text{Investimento total}}\times 100
\end{equation}

\begin{figure}[H]
\centering
\begin{tikzpicture}
\begin{axis}[width=0.86\textwidth,height=6cm,xlabel={Meses},ylabel={ROI acumulado (\%)},xmin=1,xmax=12,ymin=-25,ymax=85,grid=major,legend style={at={(0.02,0.98)},anchor=north west,font=\scriptsize}]
\addplot[color=gray!70,mark=square*] coordinates {(1,-20) (2,-16) (3,-12) (4,-6) (5,-1) (6,3) (7,8) (8,12) (9,17) (10,22) (11,26) (12,30)};
\addlegendentry{Pessimista}
\addplot[color=black,mark=*] coordinates {(1,-15) (2,-10) (3,-4) (4,2) (5,9) (6,16) (7,25) (8,34) (9,44) (10,54) (11,63) (12,72)};
\addlegendentry{Realista}
\addplot[color=gray!40,mark=triangle*] coordinates {(1,-12) (2,-6) (3,2) (4,10) (5,19) (6,30) (7,42) (8,54) (9,65) (10,73) (11,79) (12,83)};
\addlegendentry{Otimista}
\end{axis}
\end{tikzpicture}
\caption{Curvas de ROI acumulado por cenário para suporte à decisão de escala.}
\end{figure}
			
			
			%==============================================================================
			% Fim do texto
		\end{document}
